\SourcePage{195}
\SGAsection{6}{等变对应与局部项的整除性}
\textbf{6.1. 记号。} 本节中，$S$ 表示一个\Term{characteristic}{特征}为 $p$ 的域的\Term{spectrum}{谱}，$\Lambda$ 表示被某个不被 $p$ 整除的整数所\Term{annihilate}{零化}的交换\Term{noetherian-ring}{诺特环}。所考虑的 $S$-\Term{scheme}{概形}均假设\Term{separated}{分离}且\Term{finite-type}{有限型}。所考虑的 $\Lambda$-\Term{algebra}{代数}均结合、有单位且\Term{finite-type}{有限型}，但未必交换。通常，尤其在涉及“非交换的 Lefschetz--Verdier \Term{local-term}{局部项}”时，还假设这些\Term{algebra}{代数}\Term{flat}{平坦}（应用中它们将是\Term{finite-group}{有限群}的\Term{group-algebra}{群代数}）。若 $X$ 是 $S$-\Term{scheme}{概形}，$A$ 是 $\Lambda$-\Term{algebra}{代数}，以 $D({}_AX)$（分别为 $D(X_A)$）表示 $X$ 上左（分别为右）$A$-\Term{module}{模}\Term{sheaf}{层}\Term{category}{范畴}的\Term{derived-category}{导出范畴}；当 $A$ 交换时，也将 $D({}_AX)$ 写作 $D(X,A)$，并令 $D(X)=D(X,\Lambda)$。以 $D_c(-)$（分别为 $D_{ctf}(-)$）表示由具有\Term{constructible-cohomology}{可构造上同调}（分别为具有\Term{finite-tor-dimension}{有限 Tor 维数}及\Term{constructible-cohomology}{可构造上同调}）的\Term{complex}{复形}组成的\Term{full-subcategory}{满子范畴}。若 $A,B$ 是 $\Lambda$-\Term{algebra}{代数}，以 $D({}_AX_B)$ 表示 $(A,B)$-\Term{bimodule}{双模}（左 $A$、右 $B$）\Term{sheaf}{层}\Term{category}{范畴}的\Term{derived-category}{导出范畴}。最后，使用 5.1 的记号 $A^\circ,A^e,A_\natural$。

\textbf{6.2.} 第 III 讲 1、2、3 的\Term{form}{形式}体系可毫无困难地推广至“非交换”情形。下面简要说明这一推广的主要之处。

设 $X$ 是 $S$-\Term{scheme}{概形}，$A,B$ 是 $\Lambda$-\Term{algebra}{代数}。若 $L\in\operatorname{ob}D_{ctf}(X_A)$、$M\in\operatorname{ob}D({}_AX_B)$，且 $M$ 作为 $D(X_B)$ 的\Term{object}{对象}具有\Term{finite-tor-dimension}{有限 Tor 维数}及\Term{constructible-cohomology}{可构造上同调}，则 $L\overset L\otimes_A M\in\operatorname{ob}D_{ctf}(X_B)$（验证是直接的）。另一方面，由 SGA $4\frac12$ Finitude 1.6，对上述 $M$ 及 $L\in\operatorname{ob}D_{ctf}({}_AX)$，有 $R\mathcal H\!om_A(L,M)\in\operatorname{ob}D_{ctf}(X_B)$。同一出处还说明，若 $f$ 是 $S$-\Term{morphism}{态射}，则 $D_{ctf}({}_A-)$ 在运算 $(f^*,f_*,f_!,f^!)$ 下\Term{stable}{稳定}（$f^*$ 的情形显然）。

设对 $i=1,2$，$f_i:X_i\to Y_i$ 是 $S$-\Term{morphism}{态射}，
\SourcePage{196}
$f=f_1\times_S f_2:X=X_1\times_S X_2\to Y=Y_1\times_S Y_2$，$A$ 是 $\Lambda$-\Term{algebra}{代数}。对 $L_1\in\operatorname{ob}D^-(X_{1A})$、$L_2\in\operatorname{ob}D^-({}_AX_2)$，有 $D^-(Y)$ 中的 Künneth \Term{arrow}{箭头}
\begin{equation}f_{1*}L_1\overset L\otimes_{S,A}f_{2*}L_2\longrightarrow f_*(L_1\overset L\otimes_{S,A}L_2),\tag{6.2.1}\end{equation}
其中 $L_1\overset L\otimes_{S,A}L_2\overset{\mathrm{def}}=p_1^*L_1\overset L\otimes_Ap_2^*L_2$；如第 III 讲 1.6，由上述出处得知该\Term{arrow}{箭头}是\Term{isomorphism}{同构}。若 $B,C$ 是 $\Lambda$-\Term{algebra}{代数}，且假设 $B$ 或 $C$ 在 $\Lambda$ 上\Term{flat}{平坦}（以便定义 $\overset L\otimes_A:D^-({}_B-_A)\times D^-({}_A-_C)\to D^-({}_B-_C)$），则对 $L_1\in\operatorname{ob}D^-({}_BX_{1A})$、$L_2\in\operatorname{ob}D^-({}_AX_{2B})$，有 $D^-({}_BY_C)$ 中类似于 6.2.1 的\Term{isomorphism}{同构}。同样可验证，对 $M_1\in\operatorname{ob}D^-(Y_{1A})$、$M_2\in\operatorname{ob}D^-({}_AY_2)$，有 $D^-(X)$ 中的 Künneth \Term{isomorphism}{同构}
\begin{equation}f_1^!M_1\overset L\otimes_{S,A}f_2^!M_2\xrightarrow\sim f^!(M_1\overset L\otimes_{S,A}M_2).\tag{6.2.2}\end{equation}
对 $M_1\in\operatorname{ob}D^-({}_BY_{1A})$、$M_2\in\operatorname{ob}D^-({}_AY_{2C})$，也有 $D^-({}_BX_C)$ 中类似的\Term{isomorphism}{同构}。最后，设 $A_i$（$i=1,2$）是 $\Lambda$-\Term{algebra}{代数}，$A_2$ \Term{flat}{平坦}，$E_1\in\operatorname{ob}D_{ctf}({}_{A_1}X_1)$、$F_2\in\operatorname{ob}D^+({}_{A_2}X_2)$、$F_1\in\operatorname{ob}D({}_{A_1}X_{1A_2})$，且 $F_1$ 作为 $D(X_{1A_2})$ 的\Term{object}{对象}具有\Term{finite-tor-dimension}{有限 Tor 维数}及\Term{constructible-cohomology}{可构造上同调}。于是，如上所见，$R\mathcal H\!om_{A_1}(E_1,F_1)\in\operatorname{ob}D_{ctf}(X_{1A_2})$，并有 $D^+(X)$（当 $A_2$ 交换时为 $D^+(X,A_2)$）中的\Term{canonical}{典范}\Term{arrow}{箭头}
\begin{equation}R\mathcal H\!om_{A_1}(E_1,F_1)\overset L\otimes_{S,A_2}F_2\longrightarrow R\mathcal H\!om_{A_1}(p_1^*E_1,F_1\overset L\otimes_{S,A_2}F_2).\tag{6.2.3}\end{equation}
它是“\Term{inverse-image}{逆像}”\Term{arrow}{箭头}（第 III 讲 2.1.1）
\[p_1^*R\mathcal H\!om_{A_1}(E_1,F_1)\overset L\otimes_{A_2}p_2^*F_2\longrightarrow R\mathcal H\!om_{A_1}(p_1^*E_1,p_1^*F_1)\overset L\otimes_{A_2}p_2^*F_2\]
与类似于 5.4.2 的\Term{arrow}{箭头}
\[R\mathcal H\!om_{A_1}(p_1^*E_1,p_1^*F_1)\overset L\otimes_{A_2}p_2^*F_2\longrightarrow R\mathcal H\!om_{A_1}(p_1^*E_1,p_1^*F_1\overset L\otimes_{A_2}p_2^*F_2)\]
的复合；后一个\Term{arrow}{箭头}按 5.2.3 \Term{adjunction}{伴随}于由\Term{evaluation-map}{求值箭头} 5.3.6 定义的\Term{arrow}{箭头}
\[p_1^*E_1\overset L\otimes R\mathcal H\!om_{A_1}(p_1^*E_1,p_1^*F_1)\overset L\otimes_{A_2}p_2^*F_2\longrightarrow p_1^*F_1\overset L\otimes_{A_2}p_2^*F_2.\]
如第 III 讲 2.3，由 6.2.1 可得
\SourcePage{197}
6.2.3 是\Term{isomorphism}{同构}。特别地，取 $A_2=\Lambda$、$F_2=\Lambda_{X_2}$，6.2.3 给出\Term{isomorphism}{同构}
\begin{equation}p_1^*R\mathcal H\!om_{A_1}(E_1,F_1)\xrightarrow\sim R\mathcal H\!om_{A_1}(p_1^*E_1,p_1^*F_1).\tag{6.2.4}\end{equation}

\textbf{6.3.} 若 $t:T\to S$ 是 $S$-\Term{scheme}{概形}，$A$ 是 $\Lambda$-\Term{algebra}{代数}，令
\begin{equation}t^!A=K_{AT}\qquad(\in\operatorname{ob}D^+({}_AT_A)).\tag{6.3.1}\end{equation}
当 $A=\Lambda$ 时，以 $K_T$ 代替 $K_{\Lambda T}$。由 6.2.2 有 $D^+({}_AT_A)$ 中的\Term{canonical}{典范}\Term{isomorphism}{同构}
\begin{equation}K_T\overset L\otimes A\xrightarrow\sim K_{AT}.\tag{6.3.2}\end{equation}
以 $D_A$ 表示\Term{functor}{函子} $R\mathcal H\!om_A(-,K_{AT})$：它将 $D_{ctf}({}_AT)$ 映到 $D_{ctf}(T_A)$（反之亦然），而 SGA $4\frac12$ Th. Finitude 4 的证明说明，对 $L\in\operatorname{ob}D_{ctf}({}_AT)$，\Term{canonical}{典范}\Term{arrow}{箭头} $L\to D_AD_AL$ 是\Term{isomorphism}{同构}。当 $A=\Lambda$ 时，以 $D$ 代替 $D_A$。

\textbf{6.4.} 设 $A$ 是\Term{flat}{平坦} $\Lambda$-\Term{algebra}{代数}；对 $i=1,2$，设 $X_i$ 是 $S$-\Term{scheme}{概形}，$X=X_1\times_SX_2$，$p_i:X\to X_i$ 是\Term{projection}{投影}，$L_i\in\operatorname{ob}D_{ctf}({}_AX_i)$。把 6.2.3 应用于 $A_1=A_2=A$、$E_1=L_1$、$F_2=L_2$、$E_2=K_{AX_1}$，得到\Term{isomorphism}{同构}
\[D_AL_1\overset L\otimes_{S,A}L_2\xrightarrow\sim R\mathcal H\!om_A(p_1^*L_1,K_{AX_1}\overset L\otimes_{S,A}L_2).\]
与 Künneth \Term{isomorphism}{同构} 6.2.2 $K_{AX_1}\overset L\otimes_{S,A}L_2\xrightarrow\sim p_2^!L_2$ 复合，得到\Term{isomorphism}{同构}
\begin{equation}D_AL_1\overset L\otimes_{S,A}L_2\xrightarrow\sim R\mathcal H\!om_A(p_1^*L_1,p_2^!L_2),\tag{6.4.1}\end{equation}
推广了第 III 讲 3.1.1。设 $c:C\to X$ 是 $S$-\Term{morphism}{态射}，令 $c_i=p_ic$。由\Term{induction-formula}{诱导公式}（SGA 4 XVIII 3.1.12.2），有类似于第 III 讲 3.2.1 的\Term{canonical}{典范}\Term{isomorphism}{同构}
\begin{equation}c^!R\mathcal H\!om_A(p_1^*L_1,p_2^!L_2)\xrightarrow\sim R\mathcal H\!om_A(c_1^*L_1,c_2^!L_2).\tag{6.4.2}\end{equation}
与 6.4.1 复合，得到\Term{isomorphism}{同构}
\SourcePage{198}
\begin{equation}H^0(X,c_*c^!(D_AL_1\overset L\otimes_{S,A}L_2))\xrightarrow\sim\operatorname{Hom}_A(c_1^*L_1,c_2^!L_2).\tag{6.4.3}\end{equation}
$\operatorname{Hom}_A(c_1^*L_1,c_2^!L_2)$ 的元素称为从 $L_1$ 到 $L_2$、以 $c$ 为\Term{support}{支集}的\Term{cohomological-correspondence}{上同调对应}（术语与上述出处略有差异；当 $c$ \Term{proper}{固有}时两种术语一致，而实际应用中正是如此）。

给定 $S$-\Term{morphism}{态射} $f_i:X_i\to Y_i$（$i=1,2$）及第 III 讲 3.7.1 的\Term{commutative-diagram}{交换方块}，其中 $c$ \Term{proper}{固有}，则如该处一样，对 $L_i\in\operatorname{ob}D_{ctf}({}_AX_i)$，可定义类似于第 III 讲 3.7.5、3.7.6 的“\Term{direct-image}{正像}”\Term{arrow}{箭头}。

\textbf{6.5.} 设 $A$ 是\Term{flat}{平坦} $\Lambda$-\Term{algebra}{代数}，$X$ 是 $S$-\Term{scheme}{概形}，$X^2=X\times_SX$，$p_1,p_2:X^2\rightrightarrows X$ 是\Term{canonical}{典范}\Term{projection}{投影}，$\delta:X\to X^2$ 是\Term{diagonal}{对角态射}，$c:C\to X^2$ 是 $S$-\Term{morphism}{态射}，$c_i=p_ic$，$X^c=C\times_{X^2}X$ 是 $c$ 的\Term{fixed-point-scheme}{不动点概形}，因而有\Term{cartesian-square}{笛卡尔方块}
\begin{equation}\begin{tikzcd}X^c\arrow[r,"i^c"]\arrow[d,"\delta^c"']&X\arrow[d,"\delta"]\\ C\arrow[r,"c"']&X^2.\end{tikzcd}\tag{6.5.1}\end{equation}
令
\begin{equation}K_{AX,L\natural}=K_X\overset L\otimes_\Lambda(A\overset L\otimes_{A^e}A),\qquad K_{AX,\natural}=K_X\overset L\otimes_\Lambda A_\natural.\tag{6.5.2}\end{equation}
设 $L\in\operatorname{ob}D_{ctf}({}_AX)$。定义一个“\Term{local-trace}{局部迹}”\Term{homomorphism}{同态}
\begin{equation}\operatorname{Tr}_A:\delta^*R\mathcal H\!om_A(p_1^*L,p_2^!L)\longrightarrow K_{AX,L\natural},\tag{6.5.3}\end{equation}
它是以下\Term{arrow}{箭头}的复合：对 6.4.1（取 $X_1=X_2$、$L_1=L_2$）的逆施以 $\delta^*$ 所得的\Term{isomorphism}{同构}
\begin{equation}\delta^*R\mathcal H\!om_A(p_1^*L,p_2^!L)\xrightarrow\sim\delta^*(D_AL\overset L\otimes_{S,A}L),\tag{6.5.4}\end{equation}
显然的\Term{isomorphism}{同构}
\begin{equation}\delta^*(D_AL\overset L\otimes_{S,A}L)=D_AL\overset L\otimes_A L,\tag{6.5.5}\end{equation}
\SourcePage{199}
以及\Term{evaluation-map}{求值箭头} 5.3.5
\begin{equation}D_AL\overset L\otimes_A L\longrightarrow K_{AX,L\natural}.\tag{6.5.6}\end{equation}
这里借助 6.3.2 $K_X\overset L\otimes A=K_{AX}$，将 $K_{AX}\overset L\otimes_{A^e}A$ 等同于 $K_{AX,L\natural}$。由定义立即可见，当 $X=S$ 时，6.5.3 就是 5.6.1 的\Term{arrow}{箭头} $\operatorname{Tr}_A:R\mathcal H\!om_A(L,L)\to A\overset L\otimes_{A^e}A$。将对 6.5.3 施以 $i_*^c i^{c!}$ 所得的\Term{arrow}{箭头}与\Term{arrow}{箭头}
\begin{equation}\delta^*c_*R\mathcal H\!om_A(c_1^*L,c_2^!L)\longrightarrow i_*^ci^{c!}\delta^*R\mathcal H\!om_A(p_1^*L,p_2^!L)\tag{6.5.7}\end{equation}
复合，得到\Term{arrow}{箭头}
\begin{equation}\operatorname{Tr}_A:\delta^*c_*R\mathcal H\!om_A(c_1^*L,c_2^!L)\longrightarrow i_*^ci^{c!}K_{AX,L\natural}=i_*^cK_{AX^c,L\natural},\tag{6.5.8}\end{equation}
其中 6.5.7 来自 6.4.2 及\Term{base-change}{基变换}\Term{arrow}{箭头} $\delta^*c_*c^!\to i_*^ci^{c!}\delta^*$。于是得到\Term{arrow}{箭头}
\begin{equation}\operatorname{Tr}_A:\operatorname{Hom}_A(c_1^*L,c_2^!L)\longrightarrow H^0(X^c,K_{AX^c,L\natural}).\tag{6.5.9}\end{equation}
6.5.3、6.5.8、6.5.9 与由\Term{canonical}{典范}\Term{arrow}{箭头} $K_{AX,L\natural}\to K_{AX,\natural}$ \Term{derive}{导出}的\Term{arrow}{箭头}的复合，也仍记为 $\operatorname{Tr}_A$。容易验证，当 $A=\Lambda$ 时，对 $u\in\operatorname{Hom}(c_1^*L,c_2^!L)$，有
\begin{equation}\operatorname{Tr}(u)=\langle u,1\rangle,\tag{6.5.10}\end{equation}
这里使用第 III 讲 4.2.5 的记号，$1$ 表示\Term{diagonal-correspondence}{对角对应}。更确切地说，与 6.5.8 \Term{adjunction}{伴随}的\Term{arrow}{箭头} $c_*R\mathcal H\!om(c_1^*L,c_2^!L)\to\delta_*i_*^cK_{X^c}$，正是第 III 讲 4.2.4$'$（取 $X_1=X_2$、$L_1=L_2$、$d=\delta$）及\Term{diagonal-correspondence}{对角对应} $1\in\operatorname{Hom}(L,L)$ 所定义的\Term{arrow}{箭头}。

\textbf{6.6.} 设 $A$ 是\Term{flat}{平坦} $\Lambda$-\Term{algebra}{代数}，对 $i=1,2$，$X_i$ 是 $S$-\Term{scheme}{概形}，$X=X_1\times_SX_2$，$L_i\in\operatorname{ob}D_{ctf}({}_AX_i)$。如下定义\Term{pairing}{配对}
\begin{equation}R\mathcal H\!om_A(p_1^*L_1,p_2^!L_2)\overset L\otimes_\Lambda R\mathcal H\!om_A(p_2^*L_2,p_1^!L_1)\longrightarrow K_{AX,L\natural}.\tag{6.6.1}\end{equation}
先用\Term{isomorphism}{同构} 6.4.1 将左边等同于
\SourcePage{200}
\[R\mathcal H\!om_A(p_1^*L_1,p_1^*K_{AX_1})\overset L\otimes_Ap_2^*L_2\overset L\otimes_\Lambda R\mathcal H\!om_A(p_2^*L_2,p_2^*K_{AX_2})\overset L\otimes_Ap_1^*L_1,\]
再按 5.3.1 等同于
\[A\overset L\otimes_{A^e}R\mathcal H\!om_A(\ )\overset L\otimes_Ap_2^*L_2\overset L\otimes_A R\mathcal H\!om_A(\ )\overset L\otimes_Ap_1^*L_1\overset L\otimes_{A^e}A;\]
用\Term{evaluation-map}{求值箭头} 5.3.6 的\Term{tensor-product}{张量积}将此\Term{object}{对象}映到 $A\overset L\otimes_{A^e}(K_{AX_1}\overset L\otimes_AK_{AX_2})\overset L\otimes_{A^e}A=R$；最后按 5.8.2 将 $R$ 等同于 $A\overset L\otimes_{A^e}(K_{AX_1}\overset L\otimes_AK_{AX_2})$，再按 6.5.2 及第 III 讲 1.7.6 等同于 $K_{AX,L\natural}$。若 $c:C\to X$、$d:D\to X$ 是\Term{correspondence}{对应}，$e=c\times_Xd:E\to X$，则如第 III 讲 4.2，由 6.6.1 得到\Term{pairing}{配对}
\begin{equation}c_*R\mathcal H\!om_A(c_1^*L_1,c_2^!L_2)\overset L\otimes_\Lambda d_*R\mathcal H\!om_A(d_2^*L_2,d_1^!L_1)\longrightarrow e_*K_{AE,L\natural},\tag{6.6.2}\end{equation}
从而得到“\Term{cup-product}{杯积}”
\begin{equation}\langle\ ,\ \rangle_A:\operatorname{Hom}_A(c_1^*L_1,c_2^!L_2)\otimes\operatorname{Hom}_A(d_2^*L_2,d_1^!L_1)\longrightarrow H^0(E,K_{AE,L\natural}).\tag{6.6.3}\end{equation}
第 III 讲 5.2 所研究的\Term{composition-of-correspondences}{对应的复合}概念不难推广到非交换情形，并可验证，对 $u\in\operatorname{Hom}_A(c_1^*L_1,c_2^!L_2)$、$v\in\operatorname{Hom}_A(d_2^*L_2,d_1^!L_1)$，有类似于第 III 讲 5.2.9 的公式
\begin{equation}\langle u,v\rangle_A=\operatorname{Tr}_A(vu)=\operatorname{Tr}_A(uv).\tag{6.6.4}\end{equation}
Lefschetz 公式（第 III 讲 4.4）也毫无困难地推广到此情形：给定\Term{proper}{固有} $S$-\Term{morphism}{态射} $f_i:X_i\to Y_i$（$i=1,2$）及第 III 讲 4.4.0 的交换立方图，对 $L_i\in\operatorname{ob}D_{ctf}({}_AX_i)$，有类似于第 III 讲 4.4.1 的\Term{commutative-diagram}{交换方块}，其中 $K_C$（分别为 $K_D$）换为 $K_{AC,L\natural}$（分别为 $K_{AD,L\natural}$）；并有 $H^0(D,K_{AD,L\natural})$ 中类似于第 III 讲 4.5.1 的公式
\begin{equation}\langle f_*u,f_*v\rangle_A=g_*\langle u,v\rangle_A.\tag{6.6.5}\end{equation}
\Term{arrow}{箭头} $f_*c_*K_{AC,L\natural}\to d_*K_{AD,L\natural}$ 是将第 III 讲 4.4.1 的\Term{arrow}{箭头} (4) 与 $A\overset L\otimes_{A^e}A$ 作\Term{tensor-product}{张量积}而得的。
\SourcePage{201}
\textbf{6.7.} 设 $X,c$ 如 6.5，$A\to B$ 是\Term{flat}{平坦} $\Lambda$-\Term{algebra}{代数}的\Term{morphism}{态射}。关于\Term{extension-of-scalars}{标量扩张}与\Term{restriction-of-scalars}{标量限制}，\Term{arrow}{箭头} $\operatorname{Tr}_A$（6.5.3、6.5.8、6.5.9）的行为与 5.6.1、5.6.2 中的\Term{arrow}{箭头} $\operatorname{Tr}_A$ 相同。关键在于：若 $f$ 是 $S$-\Term{morphism}{态射}，则\Term{functor}{函子} $(f^*,f_*,f_!,f^!)$ 与\Term{extension-of-scalars}{标量扩张}、\Term{restriction-of-scalars}{标量限制}交换，因而研究 $\operatorname{Tr}_A$（6.5.3）关于这两种运算的行为，可归结为研究“\Term{tensor-product}{张量积}”\Term{arrow}{箭头}
\[R\mathcal H\!om_A(p_1^*L,p_1^*K_{AX})\overset L\otimes_Ap_2^*L\longrightarrow R\mathcal H\!om_A(p_1^*L,p_1^*K_{AX}\overset L\otimes_Ap_2^*L)\]
以及“\Term{evaluation-map}{求值}”\Term{arrow}{箭头}
\[R\mathcal H\!om_A(L,K_{AX})\overset L\otimes_AL\longrightarrow K_{AX,L\natural}\]
的行为。因此，方块 5.9.1、5.9.2 的\Term{commutativity}{交换性}蕴含以下方块的\Term{commutativity}{交换性}：
\begin{equation}\begin{tikzcd}[column sep=small]
\delta^*c_*R\mathcal H\!om_A(c_1^*L,c_2^!L)\arrow[r,"\operatorname{Tr}_A"]\arrow[d]&i_*^cK_{AX^c,L\natural}\arrow[d]\\
\delta^*c_*R\mathcal H\!om_B(c_1^*(B\overset L\otimes_AL),c_2^!(B\overset L\otimes_AL))\arrow[r,"\operatorname{Tr}_B"']&i_*^cK_{BX^c,L\natural}.
\end{tikzcd}\tag{6.7.1}\end{equation}
左侧竖\Term{arrow}{箭头}是\Term{extension-of-scalars}{标量扩张}，右侧竖\Term{arrow}{箭头}由\Term{canonical}{典范}\Term{arrow}{箭头} $A\overset L\otimes_{A^e}A\to B\overset L\otimes_{B^e}B$ 定义。因此，对 $u\in\operatorname{Hom}_A(c_1^*L,c_2^!L)$，有
\begin{equation}\operatorname{Tr}_B(B\overset L\otimes_Au)=\varphi\operatorname{Tr}_A(u),\tag{6.7.2}\end{equation}
其中 $\varphi:H^0(X^c,K_{AX^c,L\natural})\to H^0(X^c,K_{BX^c,L\natural})$ 是\Term{canonical-map}{典范映射}。

假设 $A$ 作为 $\Lambda$-\Term{module}{模}\Term{finite-type}{有限型}且\Term{flat}{平坦}，$B$ 作为左 $A$-\Term{module}{模}\Term{finite-type}{有限型}且\Term{flat}{平坦}，而 $A_\natural,B_\natural$ 在 $\Lambda$ 上\Term{flat}{平坦}。有 $\operatorname{Tr}_{B/A}:B_\natural\to A_\natural$，对它施以 $K_X\overset L\otimes-$（6.5.2）所得的\Term{arrow}{箭头}仍记为
\begin{equation}\operatorname{Tr}_{B/A}:K_{BX,\natural}\longrightarrow K_{AX,\natural}.\tag{6.7.3}\end{equation}
由 5.10.6、5.10.9，对 $L\in\operatorname{ob}D_{ctf}({}_BX)$，有\Term{commutative-diagram}{交换方块}
\SourcePage{202}
\begin{equation}\begin{tikzcd}
\delta^*c_*R\mathcal H\!om_B(c_1^*L,c_2^!L)\arrow[r,"\operatorname{Tr}_B"]\arrow[d]&i_*^cK_{BX^c,\natural}\arrow[d,"\operatorname{Tr}_{B/A}"]\\
\delta^*c_*R\mathcal H\!om_A(c_1^*L,c_2^!L)\arrow[r,"\operatorname{Tr}_A"']&i_*^cK_{AX^c,\natural},
\end{tikzcd}\tag{6.7.4}\end{equation}
其中左侧竖\Term{arrow}{箭头}是\Term{restriction-of-scalars}{标量限制}。所以，对 $u\in\operatorname{Hom}_B(c_1^*L,c_2^!L)$，在 $H^0(X^c,K_{AX^c,\natural})$ 中有
\begin{equation}\operatorname{Tr}_A(u)=\operatorname{Tr}_{B/A}(\operatorname{Tr}_B(u)).\tag{6.7.5}\end{equation}

\textbf{6.8.} 设 $A,B$ 是\Term{flat}{平坦} $\Lambda$-\Term{algebra}{代数}，$X,Y$ 是 $S$-\Term{scheme}{概形}，$Z=X\times_SY$，$c:C\to X^2$、$d:D\to Y^2$ 是\Term{correspondence}{对应}，$e=c\times_Sd:E=C\times_SD\to Z^2$ 是其积。因此 $Z^e=X^c\times_SY^d$。令 $R=A\otimes_\Lambda B$。如第 III 讲 5.3，定义\Term{external-tensor-product}{外张量积}
\begin{equation}\operatorname{Hom}_A(c_1^*L,c_2^!L)\otimes\operatorname{Hom}_B(d_1^*M,d_2^!M)\xrightarrow{\ \overset L\otimes_S\ }\operatorname{Hom}_R(e_1^*(L\overset L\otimes_SM),e_2^!(L\overset L\otimes_SM)).\tag{6.8.1}\end{equation}
由方块 5.12.1、5.12.3 的\Term{commutativity}{交换性}，容易\Term{pushout}{推出}：对 $u\in\operatorname{Hom}_A(c_1^*L,c_2^!L)$、$v\in\operatorname{Hom}_B(d_1^*M,d_2^!M)$，在 $H^0(Z^e,K_{RZ^e,L\natural})$ 中有
\begin{equation}\operatorname{Tr}_R(u\overset L\otimes_Sv)=\operatorname{Tr}_A(u)\operatorname{Tr}_B(v),\tag{6.8.2}\end{equation}
其中右边的乘积由\Term{canonical}{典范}\Term{arrow}{箭头}
\begin{equation}H^0(X^c,K_{AX^c,L\natural})\otimes H^0(Y^d,K_{BY^d,L\natural})\longrightarrow H^0(Z^e,K_{RZ^e,L\natural})\tag{6.8.3}\end{equation}
定义。该\Term{arrow}{箭头}来自由第 III 讲 1.7.6 及 5.9.1 中的\Term{arrow}{箭头} (4) 给出的\Term{isomorphism}{同构} $K_{AX^c,L\natural}\overset L\otimes_SK_{BY^d,L\natural}\xrightarrow\sim K_{(A\otimes B)Z^e,L\natural}$。

\textbf{6.9.} 设 $G$ 是\Term{finite-group}{有限群}。若 $X$ 是 $G$-$S$-\Term{scheme}{概形}，即带有 $G$ 的 $S$-\Term{automorphism}{自同构}\Term{action}{作用}的 $S$-\Term{scheme}{概形}\footnote{约定 $G$ 在\Term{space}{空间}上从\Term{right-action}{右作用}。}，$A$ 是 $\Lambda$-\Term{algebra}{代数}\footnote{主要关心交换 $\Lambda$-\Term{algebra}{代数}的情形。}，则以 $D(G,{}_AX)$ 表示 $X$ 上左 $G$-$A$-\Term{module}{模}\Term{sheaf}{层}\Term{category}{范畴}的\Term{derived-category}{导出范畴}（关于 $X$ 上\Term{etale-topology}{平展拓扑}的 $G$-\Term{set}{集合}\Term{sheaf}{层}、环\Term{sheaf}{层}、\Term{module}{模}\Term{sheaf}{层}等概念，参见第 X 讲 1）。
\SourcePage{203}
用下标 $c$（分别为 $ctf$）表示由如下 $G$-$A$-\Term{module}{模}\Term{complex}{复形}组成的\Term{full-subcategory}{满子范畴}：作为 $A$-\Term{module}{模}\Term{complex}{复形}，它们具有\Term{constructible-cohomology}{可构造上同调}（分别为具有\Term{finite-tor-dimension}{有限 Tor 维数}及\Term{constructible-cohomology}{可构造上同调}）。同样，若 $B$ 是 $\Lambda$-\Term{algebra}{代数}，以 $D(G,{}_AX_B)$ 表示 $X$ 上 $G$-$(A,B)$-\Term{bimodule}{双模}\Term{category}{范畴}的\Term{derived-category}{导出范畴}。若 $G$ 在 $X$ 上\Term{trivial-action}{平凡作用}，则 $X$ 上的左 $G$-$A$-\Term{module}{模}就是 $X$ 上的左 $A[G]$-\Term{module}{模}（不再带\Term{operator-group}{算子群}），所以 $D(G,{}_AX)=D({}_{A[G]}X)$；由于 $G$ \Term{finite}{有限}，还有 $D_c({}_AX)=D_c({}_{A[G]}X)$，但只有包含关系 $D_{ctf}({}_{A[G]}X)\subset D_{ctf}(G,{}_AX)$。

SGA 4 XVII、XVIII 的\Term{duality}{对偶}\Term{form}{形式}体系可毫无困难地推广到 $G$-$S$-\Term{scheme}{概形}。由于忘掉 $G$ \Term{action}{作用}的\Term{functor}{函子}\Term{faithful-functor}{忠实}，只需自然地推广“\Term{six-operations}{六种运算}”$(\overset L\otimes,R\mathcal H\!om,f^*,f_*,f_!,f^!)$ 以及该处所考虑的\Term{canonical}{典范}\Term{arrow}{箭头}（\Term{base-change}{基变换}、Künneth、\Term{duality}{对偶}等）的定义。这种推广没有困难（定义\Term{functor}{函子} $f_!,f^!$ 时，只需使用\Term{equivariant-compactification}{等变紧化}）。6.2 至 6.8 的定义与结果在 $G$-$S$-\Term{scheme}{概形}框架中的改述留给读者。若 $X_i$（$i=1,2$）是 $G$-$S$-\Term{scheme}{概形}，$c:C\to X=X_1\times_SX_2$ 是 $G$-$S$-\Term{morphism}{态射}，$L_i\in\operatorname{ob}D_{ctf}(G,{}_AX_i)$，则 $\operatorname{Hom}_A(c_1^*L_1,c_2^!L_2)$ 的元素（$\operatorname{Hom}_A$ 是 $D(G,{}_A-)$ 中的 $\operatorname{Hom}$）称为从 $L_1$ 到 $L_2$、以 $c$ 为\Term{support}{支集}的\Term{equivariant-correspondence}{等变（上同调）对应}。注意，当 $G$ 在 $X_i$（$i=1,2$）上\Term{trivial-action}{平凡作用}，且 $L_i\in\operatorname{ob}D_{ctf}({}_{A[G]}X_i)$ 时，$D(G,X)$ 中的\Term{isomorphism}{同构} 6.4.1 可以“细化”为\Term{isomorphism}{同构}
\[D_{A[G]}L_1\overset L\otimes_{S,A[G]}L_2\xrightarrow\sim R\mathcal H\!om_{A[G]}(p_1^*L_1,p_2^!L_2).\]
因此，在 6.5 类型的 $G$-\Term{equivariant}{等变}情形中，若 $G$ 在\Term{space}{空间}上\Term{trivial-action}{平凡作用}，且 $L\in\operatorname{ob}D_{ctf}({}_{A[G]}X)\subset D_{ctf}(G,{}_AX)$，则对 $u\in\operatorname{Hom}_{A[G]}(c_1^*L,c_2^!L)$，可定义\Term{local-term}{局部项} $\operatorname{Tr}_{A[G]}(u)\in H^0(X^c,K_{A[G]X^c,L\natural})$；它比仅使用 $L$ 关于 $A$ 的\Term{finite-tor-dimension}{有限 Tor 维数}所得的\Term{local-term}{局部项} $\operatorname{Tr}_A(u)\in H^0(G,X^c,K_{AX^c,L\natural})$ 更精细，其含义是显然的。

\SourcePage{204}
\textbf{6.10.} 设 $X,c$ 如 6.5，$L\in\operatorname{ob}D_{ctf}({}_{\Lambda[G]}X)$。对 $u\in\operatorname{Hom}_{\Lambda[G]}(c_1^*L,c_2^!L)$，由 6.5.9 有\Term{local-trace}{局部迹}
\begin{equation}\operatorname{Tr}_{\Lambda[G]}(u)\in H^0(X^c,K_{\Lambda[G]X^c,\natural})=H^0(X^c,K_{X^c}\otimes\Lambda[G_\natural]),\tag{6.10.1}\end{equation}
其中 $G_\natural$ 是 $G$ 的\Term{conjugacy-class}{共轭类}\Term{set}{集合}。对每个 $x\in G_\natural$，可考虑“$\operatorname{Tr}_{\Lambda[G]}(u)$ 在 $x$ 的值”$\operatorname{Tr}_{\Lambda[G]}(u)(x)\in H^0(X^c,K_{X^c})$，即 $\operatorname{Tr}_{\Lambda[G]}(u)$ 在由 $x$ 定义的\Term{projection}{投影} $H^0(X^c,K_{X^c}\otimes\Lambda[G_\natural])\to H^0(X^c,K_{X^c})$ 下的像。由 6.7.5 及 5.11.2，有
\begin{equation}\operatorname{Tr}_\Lambda(u)=|G|\operatorname{Tr}_{\Lambda[G]}(u)(e),\tag{6.10.2}\end{equation}
其中 $e$ 是 $G$ 的\Term{identity-element}{单位元}。更一般地，设 $g\in G$，以 $Z_g$ 表示 $g$ 的\Term{centralizer}{中心化子}，并以 $gu\in\operatorname{Hom}_{\Lambda[Z_g]}(c_1^*L,c_2^!L)$ 表示 $u$ 与左乘 $g:c_1^*L\to c_1^*L$ 的复合\Term{correspondence}{对应}。有 $L\in\operatorname{ob}D_{ctf}({}_{\Lambda[Z_g]}X)$，所以由 6.10.2 得到
\begin{equation}\operatorname{Tr}_\Lambda(gu)=|Z_g|\operatorname{Tr}_{\Lambda[Z_g]}(gu)(e).\tag{6.10.3}\end{equation}
用与 5.11.3 相同的推理，由 6.7.4 得到公式
\begin{equation}\operatorname{Tr}_{\Lambda[Z_g]}(gu)(e)=\operatorname{Tr}_{\Lambda[G]}(u)(g^{-1}),\tag{6.10.4}\end{equation}
因此也可将 6.10.3 写成
\begin{equation}\operatorname{Tr}_\Lambda(gu)=|Z_g|\operatorname{Tr}_{\Lambda[G]}(u)(g^{-1}).\tag{6.10.5}\end{equation}

\textbf{6.11.} 设 $f:X\to Y$ 是\Term{proper}{固有} $G$-$S$-\Term{morphism}{态射}，并给定 $G$-$S$-\Term{scheme}{概形}的 $G$-\Term{equivariant}{等变}\Term{commutative-diagram}{交换方块}
\begin{equation}\begin{tikzcd}C\arrow[r,"c"]\arrow[d]&X\times_SX\arrow[d,"f\times_Sf"]\\D\arrow[r,"d"']&Y\times_SY,\end{tikzcd}\tag{6.11.1}\end{equation}
其中 $c,d$ \Term{proper}{固有}。假设 $G$ 在 $Y,D$ 上\Term{trivial-action}{平凡作用}。以
\SourcePage{205}
\begin{equation}f^{c/d}\quad\text{（或 }f^c\text{）}:X^c\longrightarrow Y^d\tag{6.11.2}\end{equation}
表示 6.11.1 在 $c,d$ 的\Term{fixed-point-scheme}{不动点概形}上诱导的\Term{morphism}{态射}，使用 6.5.1 的记号。设 $L\in\operatorname{ob}D_{ctf}(G,X)$，$u\in\operatorname{Hom}(c_1^*L,c_2^!L)$ 是以 $c$ 为\Term{support}{支集}的 $G$-\Term{equivariant-correspondence}{等变对应}，则取\Term{direct-image}{正像}得到 $G$-\Term{equivariant-correspondence}{等变对应} $f_*u\in\operatorname{Hom}_{\Lambda[G]}(d_1^*f_*L,d_2^!f_*L)$。对 $g\in G$，以 $Z_g$ 表示 $g$ 的\Term{centralizer}{中心化子}，并以
\begin{equation}gu\in\operatorname{Hom}((gc)_1^*L,c_2^!L)\tag{6.11.3}\end{equation}
表示以
\begin{equation}gc\overset{\mathrm{def}}=(gc_1=c_1g,c_2):C\longrightarrow X\times_SX\tag{6.11.4}\end{equation}
为\Term{support}{支集}的 $Z_g$-\Term{equivariant-correspondence}{等变对应}，它定义为\Term{arrow}{箭头}
\[g^*(c_1^*L)\xrightarrow{\alpha_g}c_1^*L\xrightarrow{u}c_2^!L\]
的复合，其中 $\alpha_g$ 给出 $G$ 在 $c_1^*L$ 上的\Term{action}{作用}。显然有
\begin{equation}f_*(gu)=g(f_*u).\tag{6.11.5}\end{equation}
\begin{SGAresult}{命题}{6.12}
在 6.11 的假设下，再假设 $f_*L\in\operatorname{ob}D_{ctf}({}_{\Lambda[G]}Y)$。则对每个 $g\in G$，有
\begin{equation}f_*^{gc}\operatorname{Tr}_\Lambda(gu)=|Z_g|\operatorname{Tr}_{\Lambda[G]}(f_*u)(g^{-1}),\tag{6.12.1}\end{equation}
其中 $f_*^{gc}:H^0(X^{gc},K_{X^{gc}})\to H^0(Y^d,K_{Y^d})$ 是由\Term{adjunction}{伴随}\Term{arrow}{箭头} $f_*^{gc}K_{X^{gc}}=f_*^{gc}f^{gc!}K_{Y^d}\to K_{Y^d}$ 定义的“Gysin”\Term{morphism}{态射}。
\end{SGAresult}
\begin{proof}将 6.10.5 应用于 $f_*u$，并由 6.11.5，确有
\[\operatorname{Tr}_\Lambda(f_*(gu))=|Z_g|\operatorname{Tr}_{\Lambda[G]}(f_*u)(g^{-1}),\]
于是 6.12.1 由 Lefschetz 公式（第 III 讲 4.5.1）得到。\end{proof}

\textbf{6.13.} 设 $\ell$ 是不等于 $p$ 的素数。对每个整数 $n\ge1$，令 $\Lambda_n=\mathbf Z/\ell^n$。在 6.11 的假设下，设 $L=(L_n)$ 是
\SourcePage{206}
$D_{ctf}(G,X,\Lambda_n)$ 中\Term{object}{对象}的\Term{projective-system}{投射系}，使得当 $m\ge n$ 时 $L_m\overset L\otimes_{\Lambda_m}\Lambda_n\xrightarrow\sim L_n$；设 $u\in\operatorname{Hom}(c_1^*L,c_2^!L)$ 是\Term{equivariant-correspondence}{等变对应} $u_n\in\operatorname{Hom}(c_1^*L_n,c_2^!L_n)$ 的\Term{projective-system}{投射系}，使得对 $m\ge n$，模去\Term{isomorphism}{同构} $L_m\overset L\otimes_{\Lambda_m}\Lambda_n=L_n$ 后有 $u_m\overset L\otimes_{\Lambda_m}\Lambda_n=u_n$。由\Term{local-trace}{局部迹}与\Term{extension-of-scalars}{标量扩张}的相容性，对 $g\in G$，\Term{local-term}{局部项} $\operatorname{Tr}_{\Lambda_n}(gu_n)\in H^0(X^{gc},K_{\Lambda_n,X^{gc}})$ 定义一个元素
\[\operatorname{Tr}_{\mathbf Z_\ell}(gu)\in H^0(X^{gc},K_{\ell,X^{gc}})\overset{\mathrm{def}}=\varprojlim_n H^0(X^{gc},K_{\Lambda_n,X^{gc}}),\]
而 $f_*^{gc}\operatorname{Tr}_{\Lambda_n}(gu_n)=\operatorname{Tr}_{\Lambda_n}(gf_*u_n)$ 定义一个元素
\[f_*\operatorname{Tr}_{\mathbf Z_\ell}(gu)\in H^0(Y^d,K_{\ell,Y^d})\overset{\mathrm{def}}=\varprojlim_n H^0(Y^d,K_{\Lambda_n,Y^d}).\]
进一步假设对每个 $n$，$f_*L_n\in\operatorname{ob}D_{ctf}({}_{\Lambda_n[G]}Y)$。则\Term{local-trace}{局部迹} $\operatorname{Tr}_{\Lambda_n[G]}(f_*u_n)$ 同样定义一个元素
\[\operatorname{Tr}_{\mathbf Z_\ell[G]}(f_*u)\in H^0(Y^d,K_{\mathbf Z_\ell[G]Y^d,\natural})\overset{\mathrm{def}}=\varprojlim_n H^0(Y^d,K_{\Lambda_n[G]Y^d,\natural}).\]
由 6.12 得到
\begin{equation}f_*^{gc}\operatorname{Tr}_{\mathbf Z_\ell}(gu)=|Z_g|\operatorname{Tr}_{\mathbf Z_\ell[G]}(f_*u)(g^{-1}).\tag{6.13.1}\end{equation}

\textbf{6.14.} 设 $X$ 是 $G$-$S$-\Term{scheme}{概形}，$i:U\hookrightarrow X$ 是\Term{equivariant-open-immersion}{等变开浸入}，$c:C\to X\times_SX$ 是 $G$-$S$-\Term{scheme}{概形}的\Term{morphism}{态射}，使 $c_2(=p_2c):C\to X$ \Term{quasi-finite}{拟有限}且具有\Term{finite-tor-dimension}{有限 Tor 维数}。假设 $c_1^{-1}(U)\subset c_2^{-1}(U)$，则有\Term{commutative-diagram}{交换方块}
\begin{equation}\begin{tikzcd}[column sep=small]
U\arrow[dd,"i"']&c_1^{-1}(U)\arrow[l]\arrow[d,hook]&&\\
&c_2^{-1}(U)\arrow[d,hook]\arrow[rr]&&U\arrow[d,hook]\\
X&C\arrow[l,"c_1"]\arrow[rr,"c_2"]&&X\\
X-U\arrow[u,"j"]&c_2^{-1}(X-U)\arrow[l]\arrow[u,hook]\arrow[rr]&&X-U\arrow[u,hook].
\end{tikzcd}\tag{6.14.1}\end{equation}
\SourcePage{207}
\Term{trace-morphism}{迹态射} $\operatorname{Tr}_{c_2}:c_{2!}\Lambda_C\to\Lambda_X$ 于是定义一个\Term{filtered-category}{滤过的} $G$-\Term{equivariant-correspondence}{等变对应}
\begin{equation}\underline c\in\operatorname{Hom}(c_1^*\Lambda_X,c_2^!\Lambda_X),\tag{6.14.2}\end{equation}
其中 $\Lambda_X$ 带有由\Term{submodule}{子模} $i_!\Lambda_U$ 定义的\Term{filtration}{滤过}（即 $n\le0$ 时 $F^n\Lambda_X=\Lambda_X$，$F^1\Lambda_X=i_!\Lambda_U$，$n\ge1$ 时 $F^n\Lambda_X=0$）。其\Term{associated-graded}{相伴分次对象}由两个 $G$-\Term{equivariant-correspondence}{等变对应}组成：
\begin{equation}\begin{aligned}
\operatorname{gr}^0\underline c&=\underline c_{X-U,X}\in\operatorname{Hom}(c_1^*\Lambda_{X-U,X},c_2^!\Lambda_{X-U,X}),\\
\operatorname{gr}^1\underline c&=\underline c_{U,X}\in\operatorname{Hom}(c_1^*\Lambda_{U,X},c_2^!\Lambda_{U,X}),
\end{aligned}\tag{6.14.3}\end{equation}
其中令 $\Lambda_{X-U,X}=j_*\Lambda_{X-U}$、$\Lambda_{U,X}=i_!\Lambda_U$。有
\begin{equation}\underline c_{X-U,X}=j_*\underline c_{X-U},\tag{6.14.4}\end{equation}
其中 $\underline c_{X-U}\in\operatorname{Hom}(c_1^*\Lambda_{X-U},c_2^!\Lambda_{X-U})$ 是由 $\operatorname{Tr}_{c_2,X-U}$ 定义的\Term{correspondence}{对应}。需要指明基环时，写作 $\underline c=\underline c_\Lambda$、$\underline c_{U,X}=\underline c_{\Lambda,U,X}$ 等。

忘掉 $G$ 的\Term{action}{作用}。有\Term{local-term}{局部项}
\begin{equation}\operatorname{Tr}_\Lambda(\underline c),\quad\operatorname{Tr}_\Lambda(\underline c_{U,X}),\quad\operatorname{Tr}_\Lambda(\underline c_{X-U,X})\in H^0(X^c,K_{X^c});\tag{6.14.5}\end{equation}
由\Term{trace-additivity}{迹的可加性}（第 III 讲 4.13.1），它们满足关系
\begin{equation}\operatorname{Tr}_\Lambda(\underline c)=\operatorname{Tr}_\Lambda(\underline c_{U,X})+\operatorname{Tr}_\Lambda(\underline c_{X-U,X}).\tag{6.14.6}\end{equation}
此外，由 6.14.4，有
\begin{equation}\operatorname{Tr}_\Lambda(\underline c_{X-U,X})=j_*^c\operatorname{Tr}_\Lambda(\underline c_{X-U}),\tag{6.14.7}\end{equation}
其中 $j_*^c:H^0((X-U)^c,K_{(X-U)^c})\to H^0(X^c,K_{X^c})$ 是由 $j^c:(X-U)^c\to X^c$ 定义的\Term{direct-image}{正像}\Term{morphism}{态射}。回顾第 III 讲 4.3：若 $X$ \Term{smooth}{光滑}，$c$ 是\Term{closed-immersion}{闭浸入}且 $X^c$ \Term{finite}{有限}，则 $\operatorname{Tr}_\Lambda(\underline c)$ 是 $C$ 与\Term{diagonal}{对角线}的\Term{intersection-multiplicity}{交数} $(C.X)$ 在 $\Lambda$ 中的像；若再有 $X-U$ \Term{smooth}{光滑}，则 $\operatorname{Tr}_\Lambda(\underline c_{X-U})$ 是 $(c_2^{-1}(X-U).(X-U))$ 在 $\Lambda$ 中的像。还注意到，若 $X-U$ 在 $S$ 上\Term{proper}{固有}，且 $c|_{X-U}$ 是\Term{diagonal}{对角态射}，则 Lefschetz 公式
\SourcePage{208}
（第 III 讲 4.7）说明 $\int_{X-U}\operatorname{Tr}(\underline c_{X-U})\in\Lambda$ 是 $X-U$ 的 Euler--Poincaré \Term{euler-characteristic}{特征数}（取值于 $\Lambda$）。

\textbf{6.15.} 设给定 $G$-\Term{equivariant}{等变}图 6.11.1，其中 $f,c,d$ \Term{proper}{固有}，及 $G$-\Term{equivariant-open-immersion}{等变开浸入} $U\hookrightarrow X$。假设 $c_2$ \Term{quasi-finite}{拟有限}且具有\Term{finite-tor-dimension}{有限 Tor 维数}，$c_1^{-1}(U)\subset c_2^{-1}(U)$，$G$ 在 $Y,D$ 上\Term{trivial-action}{平凡作用}，且 $f|_U:U\to f(U)=V$ 是群为 $G$ 的\Term{galois-etale-covering}{伽罗瓦平展覆盖}。最后一个假设蕴含 $Rf_*(\Lambda_{U,X})\in\operatorname{ob}D_{ctf}({}_{\Lambda[G]}Y)$：事实上，$Rf_*(\Lambda_{U,X})$ 集中在\Term{degree}{次数} $0$，且是 $V$ 上一个\Term{local-scheme}{局部}\Term{isomorphism}{同构}于 $\Lambda[G]$ 的\Term{sheaf}{层}的\Term{extension-by-zero}{零延拓}。因此可将 6.12 应用于 $L=\Lambda_{U,X}$、$u=\underline c_{U,X}$；结合 6.14.6、6.14.7，得到
\begin{SGAresult}{命题}{6.16}
在 6.15 的假设下，对每个 $g\in G$，有
\begin{equation}f_*^{gc}\operatorname{Tr}_\Lambda(g\underline c)-(fj)_*^{gc}\operatorname{Tr}_\Lambda(g\underline c_{X-U})=|Z_g|\operatorname{Tr}_{\Lambda[G]}(f_*\underline c_{U,X})(g^{-1}),\tag{6.16.1}\end{equation}
其中 $j:X-U\hookrightarrow X$ 是\Term{inclusion}{包含映射}，$Z_g$ 是 $g$ 的\Term{centralizer}{中心化子}。
\end{SGAresult}
将 6.16 应用于\Term{correspondence}{对应} $\underline c_\ell=(\underline c_{\Lambda_n})$ 的\Term{projective-system}{投射系}，其中 $\Lambda_n=\mathbf Z/\ell^n$（$\ell$ 为不等于 $p$ 的素数）；注意到 6.14.6、6.14.7 在对 $f_*^{gc}$ 所施的映射取\Term{limit}{极限}后给出
\[f_*^{gc}\operatorname{Tr}_{\mathbf Z_\ell}(g\underline c_\ell)=f_*^{gc}\operatorname{Tr}_{\mathbf Z_\ell}(g\underline c_{\ell,U,X})+(fj)_*^{gc}\operatorname{Tr}_{\mathbf Z_\ell}(g\underline c_{\ell,X-U}),\]
结合 6.13，得到
\begin{SGAresult}{推论}{6.17}
在 6.15 的假设下，对每个 $g\in G$，有
\begin{equation}f_*^{gc}\operatorname{Tr}_{\mathbf Z_\ell}(g\underline c_\ell)-(fj)_*^{gc}\operatorname{Tr}_{\mathbf Z_\ell}(g\underline c_{\ell,X-U})=|Z_g|\operatorname{Tr}_{\mathbf Z_\ell[G]}(f_*\underline c_{\ell,U,X})(g^{-1}),\tag{6.17.1}\end{equation}
记号如 6.16。
\end{SGAresult}
由\Term{local-trace}{局部迹}与\Term{localization}{局部化}的相容性，以及 6.14 中的说明可见：当 Grothendieck 在第 XII 讲 4.5 所考虑的 $G$ 的\Term{endomorphism}{自同态} $\varphi$ 为\Term{identity-map}{恒等映射}时，6.17 以更一般的表述肯定地解决了他的整除性猜想。
\SourcePage{209}
事实上，仿照 5.13，在 $c$ 本身不 $G$-\Term{equivariant}{等变}、但 $c_2$ 是 $G$-\Term{equivariant}{等变}且 $c_1$ 对每个 $g\in G$ 满足 $c_1(xg)=c_1(x)\varphi(g)$ 的情形，不难给出类似于 6.12、6.13、6.16、6.17 的整除性命题；这里 $\varphi$ 是 $G$ 的一个给定\Term{endomorphism}{自同态}。整除性公式仍如上，只须将 $Z_g$ 换成 $g^{-1}$ 的 $\varphi$-\Term{centralizer}{中心化子}（5.13.5）。

\textbf{6.18.} 令
\begin{equation}N_G=\sum_{g\in G}g\in\Lambda[G].\tag{6.18.1}\end{equation}
设 $X$ 是 $S$-\Term{scheme}{概形}。对 $L\in\operatorname{ob}D({}_{\Lambda[G]}X)$，左乘 $N_G$ 诱导 $D(X)$ 中的\Term{arrow}{箭头}
\begin{equation}N_G:\Lambda\overset L\otimes_{\Lambda[G]}L\longrightarrow R\mathcal H\!om_{\Lambda[G]}(\Lambda,L),\tag{6.18.2}\end{equation}
其中借助自然\Term{augmentation}{增广}（把每个 $g\in G$ 映到 $1$），将 $\Lambda$ 视为 $\Lambda[G]$-\Term{algebra}{代数}。
\begin{SGAresult}{\Term{lemma}{引理}}{6.18.3}
若 $L\in\operatorname{ob}D_{ctf}({}_{\Lambda[G]}X)$，则\Term{arrow}{箭头} 6.18.2 是\Term{isomorphism}{同构}。
\end{SGAresult}
\begin{proof}事实上，通过\Term{devissage}{分解法}，可归结到 $L$ \Term{perfect-complex}{完美}，再归结到 $L=\Lambda[G]$；这时结论来自 $N_G$（6.18.1）是左 $\Lambda[G]$-\Term{module}{模} $\Lambda[G]$ 的 $G$-\Term{invariants}{不变量}\Term{module}{模}的一组基。\end{proof}

\begin{SGAresult}{命题}{6.19}
在 6.11 的情形，设 $M\in\operatorname{ob}D_{ctf}(Y)$、$v\in\operatorname{Hom}(d_1^*M,d_2^!M)$，$a:M\to f_*L$ 是 $D({}_{\Lambda[G]}Y)$ 中的\Term{arrow}{箭头}（通过\Term{augmentation}{增广} $\Lambda[G]\to\Lambda$ 将 $M$ 视为 $D({}_{\Lambda[G]}Y)$ 的\Term{object}{对象}），并使方块
\[\begin{tikzcd}d_1^*M\arrow[r,"v"]\arrow[d,"d_1^*a"']&d_2^!M\arrow[d,"d_2^!a"]\\d_1^*f_*L\arrow[r,"f_*u"']&d_2^!f_*L\end{tikzcd}\]
交换。假设 $f_*M\in\operatorname{ob}D_{ctf}({}_{\Lambda[G]}Y)$，且 $a$ 诱导\Term{isomorphism}{同构} $M\xrightarrow\sim R\mathcal H\!om_{\Lambda[G]}(\Lambda,f_*L)$。则有
\SourcePage{210}
\begin{equation}\operatorname{Tr}_\Lambda(v)=\sum_{g\in G_\natural}\operatorname{Tr}_{\Lambda[G]}(f_*u)(g^{-1}).\tag{6.19.1}\end{equation}
\end{SGAresult}
\begin{proof}由 6.18.3，$a$ 与 $N_G^{-1}$ 复合给出\Term{isomorphism}{同构} $M\xrightarrow\sim\Lambda\overset L\otimes_{\Lambda[G]}f_*L$；在此\Term{isomorphism}{同构}下，$v$ 等同于 $\Lambda\overset L\otimes_{\Lambda[G]}f_*u$。由 6.7.2，$\operatorname{Tr}_\Lambda(v)$ 就是 $\operatorname{Tr}_{\Lambda[G]}(f_*u)$ 在\Term{augmentation}{增广} $\Lambda[G]\to\Lambda$ 诱导的\Term{arrow}{箭头}下的像；而按定义，此\Term{arrow}{箭头}将每个\Term{conjugacy-class}{共轭类}映到 $1$，由此得到结论。\end{proof}

\textbf{6.20.} 设 $X_i$（$i=1,2$）是 $G$-$S$-\Term{scheme}{概形}，$X=X_1\times_SX_2$，$c:C\to X$ 是 $G$-$S$-\Term{morphism}{态射}，使得 $c_2$ \Term{quasi-finite}{拟有限}且具有\Term{finite-tor-dimension}{有限 Tor 维数}。设 $L_i\in\operatorname{ob}D^+(G,X_i)$。\Term{trace-morphism}{迹态射} $\operatorname{Tr}_{c_2}:c_{2!}c_2^*L_2\to L_2$ 经\Term{adjunction}{伴随}定义一个 $D(G,X)$ 中的\Term{arrow}{箭头}，仍记为
\begin{equation}\operatorname{Tr}_{c_2}:c_2^*L_2\longrightarrow c_2^!L_2.\tag{6.20.1}\end{equation}
借此，可给每个（$G$-\Term{equivariant}{等变}）\Term{morphism}{态射} $u\in\operatorname{Hom}(c_1^*L_1,c_2^*L_2)$ 配以一个（$G$-\Term{equivariant}{等变}）\Term{correspondence}{对应}
\begin{equation}u^!=\operatorname{Tr}_{c_2}\circ u\in\operatorname{Hom}(c_1^*L_1,c_2^!L_2).\tag{6.20.2}\end{equation}

\begin{SGAresult}{命题}{6.21}
设 $f:X\to Y$ 是\Term{proper}{固有} $G$-$S$-\Term{morphism}{态射}，并给定\Term{equivariant}{等变}\Term{commutative-diagram}{交换方块} 6.11.1，$G$ 在 $Y,D$ 上\Term{trivial-action}{平凡作用}。假设 $c,d$ \Term{proper}{固有}，$c_2,d_2$ \Term{quasi-finite}{拟有限}且具有\Term{finite-tor-dimension}{有限 Tor 维数}，并且方块
\[\begin{tikzcd}C\arrow[r,"c_2"]\arrow[d,"h"']&X\arrow[d,"f"]\\D\arrow[r,"d_2"']&Y\end{tikzcd}\]
是\Term{cartesian-square}{笛卡尔}且\Term{tor-independent}{Tor 独立}的（SGA 6 III 1.5）。设 $M\in\operatorname{ob}D_{ctf}(Y)$，$v\in\operatorname{Hom}(d_1^*M,d_2^*M)$；由此得到 $G$-\Term{equivariant}{等变}\Term{homomorphism}{同态} $h^*v\in\operatorname{Hom}(c_1^*f^*M,c_2^*f^*M)$ 以及 $G$-\Term{equivariant-correspondence}{等变对应} $(h^*v)^!\in\operatorname{Hom}(c_1^*f^*M,c_2^!f^*M)$、$f_*((h^*v)^!)\in\operatorname{Hom}(d_1^*(f_*f^*M),d_2^!(f_*f^*M))$。则有 $D({}_{\Lambda[G]}Y)$ 中\Term{arrow}{箭头}的\Term{commutative-diagram}{交换方块}（通过\Term{trivial-action}{平凡作用}将 $M$ 视为 $D({}_{\Lambda[G]}Y)$ 的\Term{object}{对象}）
\SourcePage{211}
\[\begin{tikzcd}d_1^*M\arrow[r,"v^!"]\arrow[d]&d_2^!M\arrow[d]\\d_1^*(f_*f^*M)\arrow[r,"f_*((h^*v)^!)"']&d_2^!(f_*f^*M),\end{tikzcd}\]
其中竖\Term{arrow}{箭头}由\Term{adjunction}{伴随}\Term{arrow}{箭头} $M\to f_*f^*M$ 给出。
\end{SGAresult}
\begin{proof}令 $h^*v=u$、$f^*M=L$，考虑图
\[\begin{tikzcd}[column sep=small]
d_1^*M\arrow[r,"v"]\arrow[d,"(1)"']&d_2^*M\arrow[rrr,"\operatorname{Tr}_{d_2}"]\arrow[d,"(2)"']&&&d_2^!M\arrow[d,"(3)"]\\
d_1^*f_*L\arrow[r,"(4)"']&h_*c_1^*L\arrow[r,"h_*u"']&h_*c_2^*L\arrow[r,"h_*\operatorname{Tr}_{c_2}"']&h_*c_2^!L\arrow[r,"(5)"']&d_2^!f_*L.
\end{tikzcd}\]
这里 (1)、(3) 由\Term{adjunction}{伴随}\Term{arrow}{箭头} $M\to f_*L$ 给出，(2) 由\Term{adjunction}{伴随}\Term{arrow}{箭头} $d_2^*M\to h_*h^*(d_2^*M)=h_*c_2^*L$ 给出，(4) 由\Term{base-change}{基变换}\Term{arrow}{箭头}给出，而 (5) 是\Term{canonical}{典范}\Term{isomorphism}{同构}（SGA 4 XVIII 3.1.12.3）。按定义，$v^!$ 是上行的复合。另一方面，由\Term{correspondence}{对应}的\Term{direct-image}{正像}的解释（第 III 讲 3.5 b)、3.7），$f_*u^!$ 是下行的复合。因此只需证明左侧方块 (A) 与右侧方块 (B) 交换。(A) 的\Term{commutativity}{交换性}立即来自\Term{base-change}{基变换}\Term{arrow}{箭头}的定义；至于 (B)，由\Term{tor-independent}{Tor 独立}假设，$\operatorname{Tr}_{d_2}$ 与沿 $f$ 的\Term{base-change}{基变换}相容（SGA $4\frac12$ Cycle 2.3.3），故它也交换，结论得证。\end{proof}

应用 6.19，得到：
\begin{SGAresult}{推论}{6.22}
在 6.21 的假设下，再假设 $f_*f^*M\in\operatorname{ob}D_{ctf}({}_{\Lambda[G]}Y)$，且\Term{adjunction}{伴随}\Term{arrow}{箭头} $M\to f_*f^*M$ 诱导\Term{isomorphism}{同构} $M\xrightarrow\sim R\mathcal H\!om_{\Lambda[G]}(\Lambda,f_*f^*M)$。则有
\begin{equation}\operatorname{Tr}_\Lambda(v^!)=\sum_{g\in G_\natural}\operatorname{Tr}_{\Lambda[G]}(f_*((f^*v)^!))(g^{-1}).\tag{6.22.1}\end{equation}
\end{SGAresult}

\begin{SGAresult}{命题}{6.23}
在 6.21 的假设下，再给定\Term{open-immersion}{开浸入} $i:V\hookrightarrow Y$，使得 $(fc_1)^{-1}(V)$ \Term{connected}{连通}、
\SourcePage{212}
$d_1^{-1}(V)\subset d_2^{-1}(V)$，且 $f$ 诱导一个群为 $G$ 的\Term{principal-galois-covering}{主伽罗瓦覆盖} $U=f^{-1}(V)\xrightarrow{f_V}V$。再假设 $M$ 形如 $M=i_!E$，其中 $E\in\operatorname{ob}D(V)$，并且有 $D(G,X)$ 中的\Term{isomorphism}{同构} $\alpha:f_V^*E\xrightarrow\sim\Lambda_X\otimes_\Lambda E_0$，其中 $E_0\in\operatorname{ob}D(\Lambda[G])$ 在 $\Lambda$ 上\Term{perfect-complex}{完美}。如 6.21 给定 $v\in\operatorname{Hom}(d_1^*M,d_2^*M)$，并以 $v_0\in\operatorname{Hom}_{\Lambda[G]}(E_0,E_0)$ 表示借助 $\alpha$ 由 $f^*v|_{c_1^{-1}(U)}$ 定义的\Term{homomorphism}{同态}。若 $\underline c_{U,X}$ 表示 6.14.3 所定义的\Term{correspondence}{对应}，则有（使用 6.20.2 的记号）
\begin{equation}\operatorname{Tr}_\Lambda(v^!)=\sum_{g\in G_\natural}\operatorname{Tr}_{\Lambda[G]}(f_*\underline c_{U,X})(g^{-1})\operatorname{Tr}_\Lambda(gv_0).\tag{6.23.1}\end{equation}
\end{SGAresult}
注意，由 $\operatorname{Tr}_\Lambda$ 的\Term{center}{中心}性，右边不依赖于 $\alpha$ 的选择。
\begin{proof}由\Term{projection-formula}{投影公式} $f_*f^*\Lambda_Y\otimes M\xrightarrow\sim f_*f^*M$ 及 6.15，$M$ 满足 6.22 的假设。另一方面，$\alpha$ 定义\Term{isomorphism}{同构} $f^*M\xrightarrow\sim\Lambda_{U,X}\otimes E_0$，在此\Term{isomorphism}{同构}下，$(f^*v)^!$ 等同于\Term{external-tensor-product}{外张量积} $\underline c_{U,X}\otimes v_0$。所以，由\Term{projection-formula}{投影公式}，$f_*f^*M$ 等同于 $f_*(\Lambda_{U,X})\otimes E_0$，而 $f_*((f^*v)^!)$ 等同于\Term{external-tensor-product}{外张量积} $f_*\underline c_{U,X}\otimes v_0$。对 $g\in G$，$Z_g$-\Term{equivariant-correspondence}{等变对应} $gf_*((f^*v)^!)$ 因而等同于 $gf_*\underline c_{U,X}\otimes gv_0$。由于根据 6.10.4 有
\[\operatorname{Tr}_{\Lambda[G]}(f_*((f^*v)^!))(g^{-1})=\operatorname{Tr}_{\Lambda[Z_g]}(gf_*((f^*v)^!))(e),\]
结合 6.22.1，只须证明公式
\[\operatorname{Tr}_{\Lambda[Z_g]}(gf_*\underline c_{U,X}\otimes gv_0)(e)=\operatorname{Tr}_{\Lambda[G]}(f_*\underline c_{U,X})(g^{-1})\operatorname{Tr}_\Lambda(gv_0),\]
亦即，根据 6.10.4，只须证明
\begin{equation}\operatorname{Tr}_{\Lambda[Z_g]}(gf_*\underline c_{U,X}\otimes gv_0)(e)=\operatorname{Tr}_{\Lambda[Z_g]}(gf_*\underline c_{U,X})(e)\operatorname{Tr}_\Lambda(gv_0).\tag{6.23.2}\end{equation}
将 $G$ 换成 $Z_g$ 后，可以假设 $g=e$；回到\Term{trace}{迹}的定义，可见只须证明下面的\Term{lemma}{引理}。\end{proof}

\SourcePage{213}
\begin{SGAresult}{\Term{lemma}{引理}}{6.23.3}
参见 SGA $4\frac12$ Rapport 4.7。设 $P$ 是一个 $G_V$-\Term{module}{模}，\Term{local-scheme}{局部}为\Term{finite-type}{有限型}\Term{free-module}{自由模}的\Term{direct-summand}{直和因子}，$P'=i_!P$；设 $Q$ 是一个在 $\Lambda$ 上\Term{perfect-complex}{完美}的 $\Lambda[G]$-\Term{module}{模}\Term{complex}{复形}。则通过 $G$ 的\Term{diagonal-action}{对角作用}视为 $D({}_{\Lambda[G]}V)$ 的\Term{object}{对象}的 $P\overset L\otimes_\Lambda Q$ 在 $\Lambda[G]$ 上\Term{perfect-complex}{完美}，并有\Term{commutative-diagram}{交换方块}
\[\begin{tikzcd}[column sep=small]
\delta^*R\mathcal H\!om_{\Lambda[G]}(p_1^*P',p_2^!P')\overset L\otimes R\mathcal H\!om_{\Lambda[G]}(Q,Q)\arrow[r,"(1)"]\arrow[d,"(2)"']&K_Y\arrow[d,equal]\\
\delta^*R\mathcal H\!om_{\Lambda[G]}(p_1^*P'\overset L\otimes Q,p_2^!P'\overset L\otimes Q)\arrow[r,"(3)"']&K_Y,
\end{tikzcd}\tag{$*$}\]
其中 $\delta:Y\to Y\times_SY$ 是\Term{diagonal}{对角态射}；(1) 是\Term{arrow}{箭头}
\[\delta^*R\mathcal H\!om_{\Lambda[G]}(p_1^*P',p_2^!P')\xrightarrow{6.5.3}K_{\Lambda[G]Y,\natural}\xrightarrow{x\mapsto x(e)}K_Y\]
与\Term{arrow}{箭头}
\[R\mathcal H\!om_{\Lambda[G]}(Q,Q)\xrightarrow{\operatorname{Tr}_\Lambda}\Lambda\]
的\Term{tensor-product}{张量积}；(2) 是从\Term{external-tensor-product}{外张量积}沿\Term{diagonal}{对角态射} $\Lambda[G]\to\Lambda[G]\otimes\Lambda[G]$ \Term{restriction-of-scalars}{限制}而得的\Term{arrow}{箭头}；(3) 是 6.5.3（因为 $P'\overset L\otimes Q\in\operatorname{ob}D_{ctf}({}_{\Lambda[G]}Y)$，故有定义）与 $x\mapsto x(e)$ 的复合\footnote{应当可以将 $P$ 换成一个（在 $\Lambda[G]$ 上）\Term{perfect-complex}{完美复形}，甚至将 $P'$ 换成 $D_{ctf}({}_{\Lambda[G]}Y)$ 的\Term{object}{对象}。}。
\end{SGAresult}
\begin{proof}第一个断言显然：可归结到 $P=\Lambda[G]_V$，而此时 $P\overset L\otimes_\Lambda Q$ \Term{isomorphism}{同构}于 $P$ 与带平凡 $G$ \Term{action}{作用}的 $Q$ 的\Term{tensor-product}{张量积}。对于第二个断言，先注意由 6.5.4、6.5.5 有
\[\begin{aligned}
\delta^*R\mathcal H\!om_{\Lambda[G]}(p_1^*P',p_2^!P')&\xrightarrow\sim D_{\Lambda[G]}(i_!P)\overset L\otimes_{\Lambda[G]}i_!P\\
&\xrightarrow\sim i_!(D_{\Lambda[G]}P\overset L\otimes_{\Lambda[G]}P),
\end{aligned}\]
故只须验证 $(*)$ 在 $V$ 上诱导的方块交换。因此可假设 $V=Y$。此时方块 $(*)$ 写成
\SourcePage{214}
\[\begin{tikzcd}[column sep=small]
R\mathcal H\!om_{\Lambda[G]}(P,K_Y\overset L\otimes P)\overset L\otimes R\mathcal H\!om_{\Lambda[G]}(Q,Q)\arrow[r,"(1)"]\arrow[d,"(2)"']&K_Y\arrow[d,equal]\\
R\mathcal H\!om_{\Lambda[G]}(P\overset L\otimes Q,K_Y\overset L\otimes P\overset L\otimes Q)\arrow[r,"(3)"']&K_Y,
\end{tikzcd}\tag{$**$}\]
其中 (2) 是\Term{external-tensor-product}{外张量积}，(1) 由 $\operatorname{Tr}_{\Lambda[G]}(-)(e)\otimes\operatorname{Tr}_\Lambda$ 给出（\Term{trace}{迹}取第 5 节的含义），(3) 由 $\operatorname{Tr}_{\Lambda[G]}(-)(e)$ 给出。通过\Term{descent}{下降}，可归结到 $P=\Lambda[G]_Y$；由开头的说明，此时可假设 $G$ 在 $Q$ 上\Term{trivial-action}{平凡作用}。断言于是立即由 5.12.5 得到。\end{proof}

\textbf{6.24.} 若 $\Lambda$ 被某个素数 $\ell\ne p$ 的幂所\Term{annihilate}{零化}，则可将 6.23.1 的右边改写为
\[\sum_{g\in G_\natural}\operatorname{Tr}_{\mathbf Z_\ell[G]}(f_*\underline c_{\ell,U,X})(g^{-1})\operatorname{Tr}_\Lambda(gv_0),\]
记号如 6.17。公式 6.17.1 说明：至少当 Grothendieck 在第 XII 讲 6.2 所考虑的 $G$ 的\Term{endomorphism}{自同态} $\varphi$ 为\Term{identity-map}{恒等映射}时，他在那里定义的\Term{local-term}{局部项}与 Verdier 的\Term{local-term}{局部项}一致（不过，正如已经指出的，这一\Term{restriction}{限制}并不严重）。若 $Y$ 在 $S$ 上\Term{proper}{固有}，将 6.23.1 与 $Y\to S$ 的 Lefschetz 公式（第 III 讲 4.7）结合，得到一个推广第 XII 讲 6.3 的公式；它当然也推广了 SGA $4\frac12$ Rapport 及本讲第一部分中证明的 Frobenius 的 Lefschetz 公式。

\SourcePage{215}
\SGAsection{}{参考文献}
\begin{enumerate}
\item[{[1]}] Cartan, H. 与 Eilenberg, S.：《\Term{homological-algebra}{同调代数}》（Homological Algebra），Princeton University Press，1956。
\item[{[2]}] Grothendieck, A.：《Lefschetz 公式与 $L$ 函数的有理性》（Formule de Lefschetz et rationalité des fonctions $L$），Séminaire Bourbaki 64/65，第 279 讲，收入《\Term{scheme}{概形}\Term{cohomology}{上同调}十讲》（Dix exposés sur la cohomologie des schémas），North Holland, Pub. Co.，1968。
\item[{[3]}] Illusie, L.：《\Term{cotangent-complex}{余切复形}与\Term{deformation}{形变} I、II》（Complexe cotangent et déformations I, II），Lecture Notes in Mathematics 239、283，Springer-Verlag，1971、1972。
\item[{[4]}] Langlands, R. P.：《模\Term{form}{形式}与 $\ell$ 进表示》（Modular forms and $\ell$-adic representations），收入 Modular Functions of One Variable II，Lecture Notes in Mathematics 349，Springer-Verlag，1973。
\item[{[5]}] Stallings, J.：《无\Term{center}{中心}群——Gottlieb \Term{theorem}{定理}的一个\Term{algebra}{代数}表述》（Centerless groups---an algebraic formulation of Gottlieb's theorem），Topology 4，第 129--134 页，1965。
\item[{[6]}] Verdier, J.-L.：《\Term{etale-cohomology}{平展上同调}中的 Lefschetz \Term{fixed-point}{不动点}公式》（The Lefschetz fixed point formula in étale cohomology），收入 Proceedings of a conference on local fields，Springer-Verlag，1967。
\end{enumerate}
SGA $4\frac12$：《\Term{etale-cohomology}{平展上同调}》（Cohomologie étale），P. Deligne 著，Lecture Notes in Mathematics 569，Springer-Verlag，1977。

在 SGA $4\frac12$ 中：Th. Finitude 指《$\ell$ 进\Term{cohomology}{上同调}中的\Term{finiteness-theorem}{有限性定理}》（Théorèmes de finitude en cohomologie $\ell$-adique）；Cycle 指《与一个圈相伴的\Term{cohomology-class}{上同调类}》（La classe de cohomologie associée à un cycle）；Rapport 指《迹公式报告》（Rapport sur la formule des traces）。
